\addcontentsline{toc}{chapter}{Introduction}
\chapter*{Introduction}

The field of digital manufacturing, encompassing technologies such as CNC machining and 3D printing (additive manufacturing), has become a cornerstone of modern industrial production and prototyping. The ability to translate a digital design into a physical object with micron-level precision is fundamental to sectors ranging from aerospace engineering to consumer electronics. However, despite the exponential growth in hardware capabilities, the software languages used to control these machines have remained largely static for decades. The reliance on G-code (RS-274), a low-level machine instruction language developed in the 1950s, presents a significant bottleneck in today's increasingly agile and customized manufacturing workflows.

The motivation for this project stems from the inherent limitations of the current G-code standard. While robust and universally supported, G-code lacks the structured programming features found in modern software development, such as variables, loops, conditional logic, and meaningful variable names. This deficiency forces users to rely exclusively on complex Computer-Aided Manufacturing (CAM) software to generate code, effectively treating the machine instructions as a ``black box.'' When modifications, optimizations, or experimental toolpaths are required, users face a stark choice: struggle with the opaque, error-prone editing of thousands of lines of cryptic code, or accept the limitations of their CAM software.

This project addresses a critical gap in the domain: the absence of an accessible, intermediate modelling language that sits between high-level design software and low-level machine execution. The proposed solution is a Domain-Specific Language (DSL) specifically designed for describing manufacturing operations. This DSL aims to bring the power of software engineering—modularity, reusability, and abstraction—to the physical world of machine control. By allowing users to describe \textit{what} they want to do (e.g., ``drill hole pattern'') rather than \textit{how} the motors should move, the project seeks to democratize access to advanced machine control and improve the safety and efficiency of manufacturing workflows.

The general objectives of this project are twofold. First, to perform a rigorous analysis of the problem domain, identifying the specific pain points experienced by machine operators, engineers, and researchers. Second, to design the syntax and semantics of a DSL that addresses these issues and to develop a prototype compiler capable of translating this high-level language into standard, optimized G-code.

The research methodology follows a structured Problem-Based Learning (PBL) approach. The initial phase focused on a detailed domain analysis to understand the intricacies of G-code and the limitations of existing toolchains. This involved studying the technical specifications of various machine controllers, analyzing common workflow failures, and validating the problem with potential user groups.

The structure of this report is organized as follows. \textbf{Chapter 1: Domain Analysis and Problem Definition} provides an in-depth exploration of the manufacturing domain, a detailed breakdown of the problems associated with G-code, a comparative analysis of existing methods, and the formal definition of the system requirements and constraints. \textbf{Chapter 2: Grammar Specification} details the formal language design, syntax rules, and the derivation structures of the GG-Code language. Finally, the \textbf{Conclusions} chapter summarizes the findings of the research and outlines the future direction of the project.

This introduction sets the stage for a detailed examination of why a new language for manufacturing, GG-Code, is not just a theoretical exercise, but a practical necessity for the future of flexible automation.
